\documentclass[11pt,a4paper]{article}

% Essential packages
\usepackage[T1]{fontenc}
\usepackage[utf8]{inputenc}
\usepackage{graphicx}
\usepackage{booktabs}
\usepackage{amsmath}
\usepackage{amssymb}
\usepackage{hyperref}
\usepackage{siunitx}
\usepackage{longtable}
\usepackage{multirow}
\usepackage{xcolor}
\usepackage{geometry}
\usepackage{natbib}
\usepackage{lineno}
\usepackage{subcaption}

\geometry{margin=1in}

% Hyperref setup
\hypersetup{
colorlinks=true,
linkcolor=blue,
citecolor=blue,
urlcolor=blue
}

\title{\textbf{GlycoSMILES2BAP: An Automated Pipeline for Stereochemistry-Preserving Glycan Structure Prediction with AlphaFold 3}}

\author{
\textbf{Qiang Xia}\\
Zhejiang Xinghe Tea Technology Co., Ltd.\\
Hangzhou, Zhejiang, China\\
\texttt{xiaqiang@xinghetea.com}
}

\date{}

\begin{document}

\maketitle

\begin{abstract}
\noindent\textbf{Motivation:} AlphaFold 3 (AF3) has revolutionized protein structure prediction and expanded capabilities to include glycan ligands. However, recent work identified a critical limitation: standard input formats produce stereochemically incorrect glycan structures, while the stereochemistry-preserving CCD+bondedAtomPairs (BAP) format requires impractical manual specification.

\noindent\textbf{Results:} We present GlycoSMILES2BAP, an automated pipeline that converts glycan notations to AF3-compatible format. The pipeline achieves 97.8\% epimer accuracy, 97.4\% anomeric accuracy, and 95.9\% linkage accuracy, with processing time ${<}1$ second per structure versus 30--60 minutes for manual specification. Validation against 10 literature-reported errors demonstrated 100\% correction rate, and database-scale processing of 100 GlyTouCan structures achieved 94\% success rate.

\noindent\textbf{Availability:} \url{https://github.com/xiaqiang/glycosmiles2bap}

\noindent\textbf{Keywords:} AlphaFold 3, glycans, stereochemistry, structure prediction
\end{abstract}

\linenumbers

\section{Introduction}

Glycans are essential biological molecules involved in protein folding, cell signaling, and immune recognition \citep{varki2017}. Over 50\% of human proteins undergo glycosylation \citep{helenius2004}. GlyTouCan catalogs over 200,000 unique glycan structures \citep{tiemeyer2016}.

AlphaFold 3 (AF3) represents a breakthrough in biomolecular structure prediction \citep{abramson2024}. However, \citet{huang2025} demonstrated that AF3's standard input formats fail to preserve glycan stereochemistry. SMILES format achieves only ${\sim}62\%$ accuracy, while the correct format requires manual specification taking 30--60 minutes per structure.

We present GlycoSMILES2BAP, an automated pipeline that bridges this gap by converting standard glycan notations (IUPAC-condensed, WURCS) to AF3-compatible CCD+BAP format.

\section{Methods}

\subsection{Pipeline Architecture}

GlycoSMILES2BAP operates through four sequential modules:

\begin{enumerate}
\item \textbf{Input Parser}: Converts IUPAC/WURCS to abstract syntax tree
\item \textbf{CCD Mapper}: Maps monosaccharides to Chemical Component Dictionary codes
\item \textbf{BAP Generator}: Produces bondedAtomPairs specifications
\item \textbf{Validator}: Checks stereochemistry consistency
\end{enumerate}

\subsection{CCD Mapping}

The mapper supports 28+ monosaccharide configurations. Key design decisions include:

\begin{itemize}
\item Case-insensitive matching for robust lookup
\item Anomeric position tracking (C1 for aldoses, C2 for ketoses)
\item Ring oxygen positions (O4 for pentoses, O5 for hexoses, O6 for sialic acids)
\end{itemize}

\subsection{BAP Generation}

Each glycosidic linkage is represented as an explicit atom pair specifying donor anomeric